% =====================================================================
\chapter{Formal framework}\label{ch:framework}
% =====================================================================

This chapter establishes the shared formal language. It does not invent new mathematics;
it names objects that engineering practice already uses but does not define: what a level
may ignore, under what conditions semantic preservation holds across levels, how far
comprehension should descend, and within what scope a guarantee holds.

\section{Levels and abstraction maps}\label{sec:level}

\begin{definition}[Level]\label{def:level}
A \term{level} is a six-tuple
\[
\Lambda \;=\; \bigl(X,\; \Beh,\; \Obs,\; \Env,\; \Cost,\; \Inv\bigr),
\]
where $X$ is the \term{representation space} of states admissible at that level;
$\Beh \subseteq X^{\mathbb{N}}$ is the \term{behaviour set} of admissible execution
traces; $\Obs$ is the \term{observation map} projecting the level's behaviour onto the
observable algebra of the level above; $\Env$ is a set of \term{environmental assumptions}
on which the level's conclusions depend; $\Cost : \Beh \to \mathbb{R}_{\ge 0}^{k}$ is a
\term{cost model} with $k$ components (time, space, energy, bandwidth); and
$\Inv \subseteq \Beh$ is the \term{invariant set} that must hold under every admissible
transition.
\end{definition}

The six components are not of equal weight. $X$ and $\Beh$ describe what exists, $\Obs$
describes what cannot be seen, $\Env$ describes what is depended upon, $\Cost$ describes
what it costs, and $\Inv$ describes what must not break. Chapter \ref{ch:levels} argues
that $\Obs$ and $\Inv$ decide whether a level is real: a level that cannot state its
observation map and its invariants is not a level but a directory.

\begin{definition}[Abstraction and concretisation]\label{def:galois}
Let $\Lambda_i$ and $\Lambda_{i+1}$ be adjacent levels, with the smaller index the upper
level. The \term{abstraction map} $\Obs_i : \Beh_{i+1} \to \Beh_i$ projects lower-level
behaviour onto upper-level behaviour; the \term{concretisation map}
$\Conc_i : \Beh_i \to 2^{\Beh_{i+1}}$ expands an upper-level behaviour into the set of its
lower-level implementations. If there are partial orders $\sqsubseteq_i$ and
$\sqsubseteq_{i+1}$ such that
\begin{equation}\label{eq:galois}
b \sqsubseteq_{i+1} \Conc_i(b') \quad\Longleftrightarrow\quad \Obs_i(b) \sqsubseteq_i b',
\qquad \forall b \in \Beh_{i+1},\; b' \in \Beh_i,
\end{equation}
then $(\Obs_i, \Conc_i)$ is a \term{Galois connection}, written $\Obs_i \dashv \Conc_i$.
\end{definition}

A Galois connection states precisely what abstraction discards and what it retains, and is
the central tool of the abstract-interpretation framework \cite{cousot1977}. Where an
inter-level relation can be expressed as a Galois connection, the semantic-preservation
condition becomes mechanically decidable; where it cannot, one must fall back on
simulation relations or on an explicit human argument.

\begin{definition}[Inter-level soundness]\label{def:soundness}
Let $A \in \Env_i$ be an environmental assumption declared at the upper level, and let
$\Beh_{i+1}\!\upharpoonright\! A$ denote the lower-level behaviours admissible under $A$.
Then $\Lambda_i$ is \term{sound} with respect to $\Lambda_{i+1}$ under $A$ if
\begin{equation}\label{eq:soundness}
\Obs_i\bigl(\Beh_{i+1}\!\upharpoonright\! A\bigr) \;\subseteq\; \Beh_i .
\end{equation}
\end{definition}

\begin{proposition}[Composition of soundness]\label{prop:soundness}
Let there be a chain $\Lambda_0, \Lambda_1, \dots, \Lambda_n$, with increasing index
denoting increasing concreteness, and suppose that for each $i$ there exists
$A_i \in \Env_i$ satisfying \eqref{eq:soundness}. If
\begin{enumerate}[label=(\alph*)]
\item \emph{observation compatibility}: $\Obs_i \circ \Obs_{i+1} = \Obs_i^{(i+2)}$, where
      $\Obs_i^{(i+2)}$ is the direct observation from $\Beh_{i+2}$ to $\Beh_i$; and
\item \emph{assumption closure}:
      $\Obs_i(\Beh_{i+1}\!\upharpoonright\! A_i) \subseteq \Beh_i\!\upharpoonright\! A_{i-1}$,
      that is, the assumptions satisfied below suffice to discharge those above,
\end{enumerate}
then $\Obs_0^{(n)}\bigl(\Beh_n \upharpoonright \bigcap_i A_i\bigr) \subseteq \Beh_0$.
\end{proposition}

\begin{proof}
Induction on $n$. The case $n = 1$ is \eqref{eq:soundness}. Condition (b) ensures that at
each inductive step the upper level's environmental assumption holds of the upper-level
behaviour set, so that \eqref{eq:soundness} can be applied level by level. Condition (a)
ensures that the composite observation agrees with the level-by-level observations, so
that the conclusion holds at the top level.
\end{proof}

\begin{remark}[Why composition fails]\label{rmk:why-compose-fails}
Proposition \ref{prop:soundness} is the precise form of ``how do local guarantees compose
into an end-to-end guarantee'', and equally the precise form of why they often do not:
composition requires both conditions simultaneously. In practice (a) fails when observation
maps disagree---for instance when ``the transaction committed'' denotes different
observations at the application layer and the storage layer---and (b) fails when
environmental assumptions are not closed, for instance when the upper layer assumes that
``a write is durable'' while the lower layer's durability depends on an undeclared flush
policy. Chapter \ref{ch:intent} uses these two failure modes as the skeleton of its
evidence-boundary table.
\end{remark}

\begin{remark}[Relation to the source model of the companion report]\label{rmk:source}
Definition \ref{def:level} extends the four-tuple
$M_i = (X_i, \mathrm{Beh}_i, \mathrm{Obs}_i, A_i)$ used in the companion research report
by adding $\Cost$ and $\Inv$, which Chapters \ref{ch:design} and \ref{ch:intent} require in
order to place resource guarantees and invariant enforcement inside the level object. The
companion report describes inter-level semantics by an \emph{interpretation relation}
$R_i \subseteq \mathrm{Beh}_{i+1} \times \mathrm{Beh}_i$, requiring that for every
admissible target behaviour $b_{i+1}$ of a successful transformation
$x_{i+1} = T_i(x_i)$ there exists a source behaviour $b_i$ with $R_i(b_{i+1}, b_i)$. That
is a \emph{simulation condition} and is more general than the set inclusion of
\eqref{eq:soundness}; the two coincide when $R_i$ is induced by a function, that is when
$R_i = \{(b, \Obs_i(b))\}$. This paper uses $R_i$ where a general relation is needed and
$\Obs_i$ where the abstraction is functional. The companion report stresses that, across
different languages and observation domains, $R_i$ may be a forward simulation, a backward
simulation, a bisimulation, a contextual equivalence or a refinement relation, and that
one must not assume a single relation throughout. That qualification is accepted here.
\end{remark}

\section{Oblivious sets and compression}\label{sec:oblivious}

The next two definitions are the formal centre of the paper. They turn ``what this level
need not care about'' from an engineering phrase into a computable object.

\begin{definition}[Oblivious set]\label{def:oblivious}
The \term{oblivious set} of level $\Lambda_i$ with respect to the level below is the kernel
of the abstraction map:
\begin{equation}\label{eq:oblivious}
\Obl(\Lambda_i) \;:=\; \Ker \Obs_i \;=\;
\bigl\{(b,b') \in \Beh_{i+1}^{2} \;:\; \Obs_i(b) = \Obs_i(b')\bigr\}.
\end{equation}
If $(b,b') \in \Obl(\Lambda_i)$, the level is \term{oblivious} to the difference between
$b$ and $b'$.
\end{definition}

\begin{proposition}[Basic properties]\label{prop:oblivious}
$\Obl(\Lambda_i)$ is an equivalence relation on $\Beh_{i+1}$, and the quotient satisfies
\[
\Beh_{i+1} \big/ \Obl(\Lambda_i) \;\cong\; \Obs_i(\Beh_{i+1}) .
\]
\end{proposition}

\begin{proof}
The three properties of an equivalence relation follow from those of $=$. The surjection
$\Obs_i: \Beh_{i+1} \to \Obs_i(\Beh_{i+1})$ induces a bijection
$\Beh_{i+1}/\Ker\Obs_i \to \Obs_i(\Beh_{i+1})$, by the universal property of the quotient.
\end{proof}

Proposition \ref{prop:oblivious} yields a direct consequence: \emph{abstraction is
quotienting by the kernel}. Everything a level exposes to the level above is the quotient
of lower-level behaviour by that kernel. Hence:

\begin{itemize}
\item ``what this level guarantees'' is an element of the quotient
      $\Obs_i(\Beh_{i+1})$;
\item ``what this level does not guarantee'' is what has been quotiented away, namely
      $\Obl(\Lambda_i)$;
\item ``what this level may ignore'' is every lower-level detail compatible with
      $\Obl(\Lambda_i)$.
\end{itemize}

The third item has the deepest consequence: it makes ignorance a property that can be
declared, checked and held to account. A module that has never declared its oblivious set
has an incomplete contract.

\begin{remark}[A criterion for pattern realism]\label{rmk:dennett}
Proposition \ref{prop:oblivious} is isomorphic to an independent philosophical criterion.
Dennett's criterion for a \emph{real pattern} states that a pattern is real exactly when
describing the system in terms of it is more economical than describing the raw data
\cite{dennett1991}. Under that criterion, whether a level exists is not an ontological
question but a question of compression ratio---precisely the quantity measured by
Definition \ref{def:compression}.

Three consequences follow, and are used in later chapters.
\begin{enumerate}[label=(\alph*)]
\item Levels are neither discovered natural kinds nor arbitrary conventions; they are fixed
      points of a compression operator.
\item The convergence thesis of Chapter \ref{ch:convergence} needs no separate
      philosophical defence: if two systems compress independently under similar
      constraints, that their fixed points are similar follows from properties of the
      compression operator rather than from an additional assumption about the world.
\item The criterion is quantitative, so whether a given level is real is measurable rather
      than arguable.
\end{enumerate}
\end{remark}

\begin{definition}[Compression ratio]\label{def:compression}
Let $\mu$ be a reference measure on $\Beh_{i+1}$ and let $\Entropy$ denote the
corresponding entropy functional. The \term{compression ratio} of level $\Lambda_i$ is the
information annihilated:
\begin{equation}\label{eq:compression}
\kappa_i \;=\; \Entropy_{\mu}\bigl(\Beh_{i+1}\bigr) \;-\;
\Entropy_{\Obs_i \# \mu}\bigl(\Obs_i(\Beh_{i+1})\bigr),
\end{equation}
that is, the information that cannot be recovered under $(\Obs_i, \Conc_i)$. The wording
is borrowed from information theory \cite{shannon1948} and algorithmic information theory
\cite{kolmogorov1965}.
\end{definition}

Compression is not itself the objective: a system that guarantees nothing attains
arbitrarily large $\kappa$. The objective is to maximise $\kappa$ \emph{subject to
preserving task-relevant predictive power}, which is the subject of
\S\ref{sec:comprehension-formal}.

\section{Architecture as a typed relational graph}\label{sec:relations}

Describing a system by levels alone is insufficient. Chapter \ref{ch:convergence} argues
that cross-system similarity resides mainly in \emph{relations} rather than in components,
so the relations themselves must be typed.

\begin{definition}[Architecture relational graph]\label{def:relgraph}
The \term{architecture relational graph} of a system is a typed directed multigraph
\begin{equation}\label{eq:relgraph}
\mathcal{G} \;=\; \bigl(V,\; E_{\mathrm{abs}} \sqcup E_{\mathrm{dep}} \sqcup
E_{\mathrm{cmp}} \sqcup E_{\mathrm{ctl}} \sqcup E_{\mathrm{ctr}}\bigr),
\end{equation}
where $V$ is a set of components and the five edge classes are:
\begin{enumerate}[label=(\alph*)]
\item $E_{\mathrm{abs}}$, \term{abstraction}: $u \to v$ when several details of $v$ are
      expressed uniformly by $u$; annotated with which semantics of $v$ are retained and
      which are hidden.
\item $E_{\mathrm{dep}}$, \term{dependency}: $u \to v$ when $u$ requires $v$ to discharge
      its own duty; annotated with whether the dependency is on an abstract interface or on
      a concrete implementation.
\item $E_{\mathrm{cmp}}$, \term{composition}: $u \to v$ when $u$ is constituted by $v$;
      annotated with ownership versus mere use.
\item $E_{\mathrm{ctl}}$, \term{control}: $u \to v$ when execution passes from $u$ to $v$;
      annotated with synchronous call, asynchronous message, event callback or data-driven
      dispatch, and with who schedules.
\item $E_{\mathrm{ctr}}$, \term{contract}: $u \to v$ when $u$ promises something to $v$ and
      requires something of $v$; annotated with preconditions, postconditions, invariants
      and failure semantics.
\end{enumerate}
\end{definition}

Keeping the five classes separate is not formalism for its own sake. Consider a single
call graph: it would compress four distinct facts into one edge class---that a thread pool
composes threads and queues, that it depends on synchronisation primitives, that it is
invoked through a submission interface, and that it promises to accept no new tasks after
shutdown---and would lose all structural information in doing so. The seam detectors of
Chapter \ref{ch:comprehension} are, in effect, procedures for recovering this five-way
distinction.

\begin{definition}[Abstract skeleton]\label{def:skeleton}
A \term{skeleton map} $\phi$ replaces each component by its role label in some duty
classification and retains the type of each relation while discarding its instance. The
image $\phi(\mathcal{G})$ is the \term{abstract skeleton} of $\mathcal{G}$.
\end{definition}

The central proposition of Chapter \ref{ch:convergence} is stated with $\phi$:
$\phi(\mathcal{G}_1) \approx \phi(\mathcal{G}_2)$ can hold even when
$\mathcal{G}_1$ and $\mathcal{G}_2$ have nothing in common at the source level.

\section{Comprehension as task-constrained lossy compression}\label{sec:comprehension-formal}

\begin{definition}[Task]\label{def:task}
A \term{task} is a pair $T = (Q, d)$, where $Q$ is a family of \term{queries}
$q : \Beh \to \mathcal{A}_q$ mapping system behaviour to observable answers, and $d$ is a
loss functional on answers.
\end{definition}

Typical queries are: if I change this, which behaviours change? under what conditions does
this path fail? who enforces this invariant? which modules must change to add a new entity
type?

\begin{definition}[$T$-sufficiency]\label{def:sufficient}
Let $M$ be a model of system $S$ at some level of abstraction. $M$ is
\term{$T$-sufficient} for $S$ if
\begin{equation}\label{eq:sufficient}
\forall q \in Q:\quad q\bigl(\Beh(M)\bigr) = q\bigl(\Beh(S)\bigr),
\end{equation}
or, in the approximate form,
$\mathbb{E}_{q \sim Q}\bigl[d\bigl(q(\Beh(M)), q(\Beh(S))\bigr)\bigr] \le D_T$.
\end{definition}

\begin{definition}[Description length and the rate--distortion objective]\label{def:mdl}
The definition rests on the minimum description length principle
\cite{rissanen1978,grunwald2006,grunwald2019}. Let $\DL(M)$ be the description length of
$M$ under a fixed encoding, and let
$L_T(M) = \mathbb{E}[d(q(\Beh(M)), q(\Beh(S)))]$ be the task loss. Comprehending $S$ means
solving
\begin{equation}\label{eq:ratedistortion}
M^{\star} \;=\; \arg\min_{M} \; \DL(M)
\qquad \text{s.t.} \qquad L_T(M) \le D_T ,
\end{equation}
or equivalently the Lagrangian form $\min_M \; \DL(M) + \lambda\, L_T(M)$.
\end{definition}

\begin{proposition}[Stopping rule for descent]\label{prop:stopping}
Parameterise the model family by descent depth, $M_0 \succ M_1 \succ \cdots$, with
increasing index meaning increasing retention of lower-level detail, and suppose
$\DL(M_{k+1}) \ge \DL(M_k)$ and $L_T(M_{k+1}) \le L_T(M_k)$. Then the optimal depth
$k^{\star}$ satisfies the stationarity condition
\begin{equation}\label{eq:stopping}
\left.\frac{\partial L_T}{\partial \DL}\right|_{k^{\star}} \;=\; -\frac{1}{\lambda} ,
\end{equation}
and once $\partial L_T / \partial \DL = 0$ beyond some depth, further descent is strictly
dominated.
\end{proposition}

\begin{proof}
The Lagrangian is $\mathcal{L}(k) = \DL(M_k) + \lambda L_T(M_k)$. First-order optimality
gives $\mathrm{d}\DL/\mathrm{d}k + \lambda\, \mathrm{d}L_T/\mathrm{d}k = 0$, that is
$\mathrm{d}L_T/\mathrm{d}\DL = -1/\lambda$. If
$\mathrm{d}L_T/\mathrm{d}\DL = 0$ on an interval while $\mathrm{d}\DL/\mathrm{d}k > 0$,
then $\mathrm{d}\mathcal{L}/\mathrm{d}k = \mathrm{d}\DL/\mathrm{d}k > 0$ and the objective
increases monotonically, so further descent is strictly dominated.
\end{proof}

Proposition \ref{prop:stopping} yields a rule usable in practice.

\begin{principle}[Descent depth]\label{pr:descent}
To make a change correctly it suffices to descend until the invariant governing that change
becomes visible, and then to stop. Further descent increases description length without
reducing task loss.
\end{principle}

The rule explains both why reading all the source is not comprehension and why reading only
an architecture diagram is not comprehension either: the former passes the stationary point
into the dominated region, the latter has not yet reached $T$-sufficiency.

\section{Contracts, traces and assurance}\label{sec:contract}

\begin{definition}[Behaviour contract]\label{def:contract}
A \term{behaviour contract} is an eight-tuple
\begin{equation}\label{eq:contract}
S \;=\; \bigl(P,\; Q,\; E^{+},\; E^{-},\; R,\; I,\; F,\; A\bigr),
\end{equation}
where $P$ is a precondition; $Q$ a postcondition; $E^{+}$ the set of permitted side
effects; $E^{-}$ the set of forbidden side effects; $R$ a resource budget (time, space,
bandwidth, energy); $I$ the invariants that must hold throughout execution; $F$ the
failure, cancellation, retry and recovery semantics; and $A$ the environmental assumptions
(hardware, operating system, load, external services).
\end{definition}

Equation \eqref{eq:contract} is not claimed to be a recognised specification format. Its
purpose is to state that \emph{correctness is not compressed into a return value but
defined as a set of constraints on the whole execution}. That is the starting point of
Chapter \ref{ch:intent}.

\begin{definition}[Traces and the admissible set]\label{def:traces}
Let $B$ be an executable artefact and $E$ a runtime environment (processor, instruction
set, operating system, loader, configuration, external services and devices). Let
$\Tr(B,E)$ be the set of \emph{observable execution traces} $B$ may produce in $E$,
including input and output, system calls, state mutations, I/O, concurrency events and
failure behaviour, and let $\Allow(S)$ be the traces permitted by contract $S$. Then $B$
\term{satisfies} $S$ in $E$ if
\begin{equation}\label{eq:inclusion}
\Tr(B,E) \;\subseteq\; \Allow(S).
\end{equation}
\end{definition}

Inclusion \eqref{eq:inclusion} is far stronger than output correctness: it constrains
\emph{all possible behaviours}, including side effects that were never intended.

\begin{definition}[Assurance and its condition]\label{def:assurance}
An \term{assurance} is a four-tuple (property assured, object, premises, evidence). A
property $\varphi$ holds within scope $\Sigma$ exactly when
\begin{equation}\label{eq:assurance}
\begin{aligned}
&\text{environmental assumptions hold} \\
&\land\; \text{the contract is explicit and authorised} \\
&\land\; \text{the verifier accepts the implementation and its proof} \\
&\land\; \text{the transformations preserve the required semantics} \\
&\land\; \text{the execution boundary is effective} \\
&\land\; \text{the object executed is the object verified} \\
&\qquad\Longrightarrow\; \varphi \text{ holds within the declared scope } \Sigma .
\end{aligned}
\end{equation}
\end{definition}

The last conjunct---called here the \term{fourth pillar}---is frequently omitted, but it
is not a detail: even if a verifier correctly accepts a binary, the preceding conclusion
does not transfer if the deployment system actually runs a different file, or if loading
changes a critical premise. Its formal expression is an identity assertion:
\begin{equation}\label{eq:binding}
\mathrm{digest}\bigl(B_{\text{deployed}}\bigr) \;=\;
\mathrm{digest}\bigl(B_{\text{verified}}\bigr)
\quad \land \quad
E_{\text{deployed}} \models A_{\text{claimed}} .
\end{equation}
Equation \eqref{eq:binding} is why content addressing, reproducible builds and load-chain
integrity belong to the theory rather than to deployment housekeeping.

\section{Two modes of trust generation}\label{sec:trust}

\begin{definition}[Trusted computing base]\label{def:tcb}
For an assurance regime, the \term{trusted computing base} $\TCB$ is the set of components
that must be trusted unconditionally for the assurance to hold: the proof checker and its
execution, the logic and semantic definitions in use, the identity between the program a
proof is about and the program that executes, the compilation, linking, loading and
runtime mechanisms, and the assumed hardware and environmental conditions.
\end{definition}

\term{TCB minimisation} is the precise restatement of the intuition that a system must not
verify itself. Its exact content is this: a mechanical system does not \emph{betray}, but
it can be defective, can fail, and can embody an incomplete model. The basis of trust is
therefore not that the mechanism is error-free, but that its decision does not depend on
the generator's will. The two are not the same.

\begin{definition}[Mechanical and consensus assurance]\label{def:twotrust}
\term{Mechanical assurance} is provided by a single deterministic, exhaustively inspectable
decision procedure whose objects are finite, static and unambiguous symbolic objects, such
as proof terms, typing derivations or constraint systems. \term{Consensus assurance} is
provided by the adversarial scrutiny of multiple independent participants, and its objects
are high-dimensional, dynamic and partly self-referential judgements, such as whether a
specification faithfully expresses an intent.
\end{definition}

The two modes are complementary rather than competing. Table \ref{tab:twotrust} states the
difference.

\begin{table}[htbp]
\centering
\caption{The two modes of trust generation}\label{tab:twotrust}
\begin{tabular}{L{2.7cm}L{5.3cm}L{5.3cm}}
\toprule
\textbf{Dimension} & \textbf{Mechanical assurance} & \textbf{Consensus assurance} \\
\midrule
Generation & A minimal decision procedure examines symbolic objects &
  Independent parties scrutinise adversarially and sign \\
Object & Formalised properties & Judgements that cannot be formalised \\
Strength & Absolute under the declared premises & Statistically reliable, not absolute \\
Failure mode & Defective decision procedure; lapsed premise &
  Collusion, capture, \emph{cognitive isomorphism} \\
Failure shape & Discrete (defective or not) & Continuous (higher independence, lower failure) \\
Cost & Proportional to the complexity of the property & Proportional to participant count and independence \\
Disciplinary analogue & Checking a mathematical proof & Peer review, a jury, the judiciary \\
\bottomrule
\end{tabular}
\end{table}

The load-bearing row is failure shape. Mechanical assurance fails discretely---not because
it is more reliable, but because its object is a finite static object and the check can be
exhausted. Consensus assurance fails continuously, because its object is intentional and
can only be approached by a network. \S\ref{sec:consensus} gives the precise form of that
continuity and proves a counter-intuitive consequence: the governing variable of consensus
reliability is not the number of participants but the independence of their errors.

\section{Contract operads and composition}\label{sec:operad}

Proposition \ref{prop:soundness} gives sufficient conditions for composition, but it lists
them one by one. This section supplies the algebraic structure of contracts, so that the
composition conditions become well-definedness conditions on that structure rather than a
checklist.

\begin{definition}[Contract operad]\label{def:operad}
A \term{contract operad} $\mathcal{O}$ consists of:
\begin{enumerate}[label=(\roman*)]
\item a set of \term{colours} $\mathrm{Col}$, whose elements are behaviour contracts
      (Definition \ref{def:contract});
\item for each $n \ge 0$ and colours $S_0, S_1, \dots, S_n$, a set of \term{operations}
      $\mathcal{O}(S_1, \dots, S_n;\, S_0)$ whose elements are \emph{composites taking
      $S_1,\dots,S_n$ as subcomponent contracts and satisfying $S_0$};
\item \term{substitution}: from $f \in \mathcal{O}(S_1,\dots,S_n;\,S_0)$ and
      $g_i \in \mathcal{O}(T_{i,1},\dots,T_{i,m_i};\,S_i)$ for $i = 1..n$,
      \[
      f \circ (g_1, \dots, g_n) \;\in\;
      \mathcal{O}\bigl(T_{1,1},\dots,T_{1,m_1},\dots,T_{n,1},\dots,T_{n,m_n};\,S_0\bigr);
      \]
\item \term{units} $\mathrm{id}_S \in \mathcal{O}(S;\,S)$.
\end{enumerate}
Substitution is required to be associative and unital. Informally, a colour is what a part
promises, an operation is how parts are assembled, and substitution refines a part without
altering the promise of the whole. The definition is the standard one for an operad,
coloured by contracts \cite{leinster2004}.
\end{definition}

\begin{definition}[Interface graph and freeness]\label{def:interface-graph}
Let $f \in \mathcal{O}(S_1,\dots,S_n;\,S_0)$. Take the subcomponents as vertices and draw a
directed edge $u \to v$ when the guarantee of $u$ is required by an assumption of $v$. The
resulting digraph is the \term{interface graph} of $f$. If it is acyclic, $f$ is said to be
\term{freely generated} by its interface graph.
\end{definition}

\begin{proposition}[Structural form of the composition conditions]\label{prop:operad-compose}
Conditions (a) and (b) of Proposition \ref{prop:soundness} are respectively equivalent to
two structural conditions on the contract operad:
\begin{enumerate}[label=(\alph*)]
\item \term{Well-typed colouring}: observation compatibility (condition (a) of
      Proposition \ref{prop:soundness}) holds if and only if substitution is well defined
      on colours, that is, the composite never demands of a subcomponent an observation its
      colour does not declare;
\item \term{Environmental closure}: assumption closure (condition (b) of
      Proposition \ref{prop:soundness}) holds if and only if the interface graph is
      acyclic, that is, the composite can be built uniquely bottom-up from its
      subcomponents.
\end{enumerate}
\end{proposition}

\begin{proof}
For (a): if the observation maps are compatible, the projection of each subcomponent onto
the interface agrees with the composite's observation, so substitution introduces no new
colour requirement; conversely, if substitution is well typed on colours, every observable
of the composite is traceable to an observation declared by some subcomponent colour, which
is observation compatibility.

For (b): if the interface graph is acyclic, a topological order exists and substitution
along that order yields the unique construction, with assumption closure verified level by
level. Conversely, if there is a cycle $u_1 \to u_2 \to \cdots \to u_k \to u_1$, then the
assumption of $u_1$ depends on the guarantee of $u_k$, whose guarantee depends (around the
cycle) on the assumption of $u_1$; the two sides of the substitution define each other and
the composite has no well-founded construction. To give it meaning one must supply a
fixed-point semantics, that is, a separate existence-and-uniqueness condition for that
cycle.
\end{proof}

Part (b) accounts for two of the three composition failures in \S\ref{sec:composition}, and
explains why two established practices must introduce an additional hypothesis:

\begin{itemize}
\item contract composition for dynamical systems must assume that the interconnection is
      \emph{well posed}, which is exactly the fixed-point semantics supplied for a feedback
      cycle;
\item software contracts fail on network-retry loops because retry turns an acyclic call
      graph into a cyclic interface graph, and the colouring is not re-verified once the
      cycle appears.
\end{itemize}

\begin{remark}[Feedback is a structural change]\label{rmk:feedback}
Proposition \ref{prop:operad-compose} reduces composability to acyclicity of the interface
graph. This has a non-trivial consequence: \emph{introducing feedback is not a local
change}. Any edit that turns an acyclic interface graph into a cyclic one---retry, cache
write-back, bidirectional synchronisation, an event loop---invalidates the existing
composition argument, regardless of whether the components remain individually correct.
This agrees with engineering experience: such edits are a high-incidence region for
regressions.
\end{remark}

\section{Notation}\label{sec:notation}

\begin{longtable}{L{3.2cm}L{11.3cm}}
\caption{Notation used throughout.}\label{tab:notation}\\
\toprule
\textbf{Symbol} & \textbf{Meaning} \\
\midrule
\endfirsthead
\toprule
\textbf{Symbol} & \textbf{Meaning} \\
\midrule
\endhead
\bottomrule
\endlastfoot
$\Lambda_i$ & level: the six-tuple (Definition \ref{def:level}) \\
$\Obs_i \dashv \Conc_i$ & Galois connection between levels (Definition \ref{def:galois}) \\
$\Obl(\Lambda_i) = \Ker \Obs_i$ & oblivious set (Definition \ref{def:oblivious}) \\
$\kappa_i$ & compression ratio (Definition \ref{def:compression}) \\
$\mathcal{G}$ & architecture relational graph (Definition \ref{def:relgraph}) \\
$\phi(\mathcal{G})$ & abstract skeleton (Definition \ref{def:skeleton}) \\
$T = (Q,d)$ & task (Definition \ref{def:task}) \\
$\DL(M)$ & description length (Definition \ref{def:mdl}) \\
$L_T(M)$ & task loss (Definition \ref{def:mdl}) \\
$S$ & behaviour contract: the eight-tuple (Definition \ref{def:contract}) \\
$\Tr(B,E)$, $\Allow(S)$ & observable traces, permitted traces
  (Definition \ref{def:traces}) \\
$\mathfrak{A}$ & assurance (Definition \ref{def:assurance}) \\
$\TCB$ & trusted computing base (Definition \ref{def:tcb}) \\
$\mathcal{O}$, $\mathrm{Col}$ & contract operad and its colours (Definition \ref{def:operad}) \\
$R_i$ & inter-level interpretation relation (Remark \ref{rmk:source}) \\
$\rho(c)$ & read-set of change $c$ (Definition \ref{def:readset}) \\
$\mathcal{K}_6$ & the six functions (Definition \ref{def:six}) \\
$D$ & technical debt (Definition \ref{def:debt}) \\
$C(t)$, $K$ & governable complexity, comprehension capacity
  (Definition \ref{def:governable}) \\
$\mathcal{J}$ & observation operator (Definition \ref{def:observability}) \\
$\operatorname{rank}_{\mathcal{T}}\mathcal{J}$ & observability rank on the task subspace
  (Definition \ref{def:observability}) \\
$q$ & correlated error mode (\S\ref{sec:consensus}) \\
\end{longtable}
